\section{Summary}
\label{sec:conclusion}

Box transpose fixes a tensor's layout in place. We tested it on \resCases{}
sampled shapes at ranks \resRankLo{} to \resRankHi{}. We also tested it on
\hardCases{} permutations picked at random. It found a working plan and gave
the right answer every time. It used \resBoxAux{} bytes of scratch space. A
copy would have used a second tensor of \resTensorMB{}\,MB.

When a permutation leaves a group of axes at the back alone, we are the fastest
in-place method we measured. We beat ITTPD, Gomez-Luna and Gustavson on
\resWinIttpd{}, \resWinGl{} and \resWinGv{} of their cases.

When the swap is also one small square block, we beat PyTorch's own copy. We
are \tierBeatsX{}$\times$ faster. The reason is simple. We skip the cells that
are already in the right place. A copy cannot.

We work out both tests from the shape, before touching any data. No number in
the cost model was measured on a GPU.

When the tests fail, our method is the slowest one we compared. That happens
when the permutation moves the innermost axis. We are about
$\hardSmallSlowdown\times$ slower there. We report it instead of routing around
it.
